\begin{ejercicio}
    % ID: 3
        Mostrar que la homotopía es una relación de equivalencia en el conjunto de funciones continuas de $X$ en $Y$, denotado por $\text{Map}(X,Y)$.
    \end{ejercicio}

\begin{demostracion}
(1) Sea $f: X\to Y$ una función continua. Hagamos 
\begin{align*}
    H:X\times I&\to Y\\
    (x,t)&\mapsto f(x).
\end{align*}
Se sigue que $H$ es continua y cumple que $H(x,0)=H(x,1)=f(x)$ para toda $x\in X$. Por lo tanto, $f\simeq f$. (2) Supongamos $f\simeq g,$ entonces existe $H:X\times I\to Y$ continua tal que 

$$H(x,0)=f(x)\text{\; y\; } H(x,1)=g(x).$$
Hagamos $H_1:X\times I\to Y$ tal que $H_1(x,t)=H(x,1-t)$. Se sigue que $H_1$ es continua y cumple que 
\begin{align*}
    H_1(x,0)&=H(x,1)=g(x),\\
    H_1(x,1)&=H(x,0)=f(x)
\end{align*}
para toda $X\in X$. Por lo tanto, $H_1$ es testigo de que $g\simeq f$. (3) Supongamos $f\simeq g$ y $g\simeq h$. Entonces, existen $H_1,H_2:I\times X\to Y$ continuas tales que
\begin{align*}
    H_1(x,0)&=f(x),\\
    H_1(x,1)&=g(x)=H_2(x,0)\text{\; y }\\
    H_2(x,1)&=h(x).
\end{align*}

Definamos $H:I\times X\to Y$ tal que
\begin{equation*}
    H(x,t)=\begin{cases}
        H_1(x,2t), & \text{ si } t\in[0,\nicefrac{1}{2}];\\[0.25ex]
        H_2(x,2t-1) & \text{ si } t\in [\nicefrac{1}{2},1].
    \end{cases}
\end{equation*}
Notemos que $H_1(x,1)=H_2(x,1)=g(x)$, para toda $x\in X$, por el lema del pegado, $H$ es continua. Se sigue que 
\begin{align*}
    H(x,0)&=H_1(x,0)=f(x),\\
    H(x,1)&=H_2(x,1)=h(x)
\end{align*}
para toda $x\in X$. Se sigue que $H$ es testigo de que $f\simeq h$.
\end{demostracion}
